\documentclass[12pt]{article}
\usepackage{graphicx} % Required for inserting images
\usepackage[margin=0.5in]{geometry}
\usepackage{amsmath, amssymb}
\usepackage{mathrsfs}


\usepackage{soul}


% Dr. David Brown Commands
\newcommand{\ds}{\displaystyle}
\newcommand{\R}{\mathbb{R}}
\newcommand{\N}{\mathbb{N}}
\newcommand{\Q}{\mathbb{Q}}
\newcommand{\C}{\mathbb{C}}


% Commands to get script r
\usepackage{calligra}
\DeclareMathAlphabet{\mathcalligra}{T1}{calligra}{m}{n}
\DeclareFontShape{T1}{calligra}{m}{n}{<->s*[2.2]callig15}{}
\newcommand{\scripty}[1]{\ensuremath{\mathcalligra{#1}}}

% Unit commands
\newcommand{\degree}{\text{°}}
\newcommand{\unitSpeed}{\frac{\text{m}}{\text{s}}}
\newcommand{\unitAcceleration}{\frac{\text{m}}{\text{s}^2}}

% Constant commands
\newcommand{\avogadro}{6.022\times10^{23}}

\newcommand{\boltzmannConstK}{1.381\times10^{-23}\frac{\text{J}}{\text{K}}}
\newcommand{\boltzmannConstInEV}{8.617\times10^{-5}\frac{\text{eV}}{\text{K}}}
\newcommand{\planckConst}{6.626\times10^{-34}\text{J}\cdot\text{s}}

\newcommand{\atmP}{1.01325\times10^5\,\text{Pa}}

% Commands to easily write partial derivatives
\newcommand{\partDeriv}[2]{\frac{\partial #1}{\partial #2}}
\newcommand{\doublePartDeriv}[2]{\frac{\partial^2 #1}{\partial^2 #2}}


\title{Problem 6.26}
\author{Gabriel Decker}


\begin{document}



\maketitle
\begin{flushleft}

In this problem, we deal with the low-temperature limit $(kT\ll\epsilon)$ for the rotational partition function. Because of this limit, we can accurately approximate results by just keeping the largest $T$-dependent term at each part of the calculation.
The partition function $Z$ is given by the following equation.
\begin{equation}
    Z_\text{rot}=\sum_{j=0}^\infty(2j+1)e^{-j(j+1)\epsilon/kT}
\end{equation}\\~\\

We'll start by plugging in a few $j$ values.
$$Z_\text{rot}=\sum_{j=0}^\infty(2j+1)e^{-j(j+1)\epsilon/kT}$$
$$\implies Z_\text{rot}\approx\sum_{j=0}^2(2j+1)e^{-j(j+1)\epsilon/kT}$$
$$\implies Z_\text{rot}\approx e^{0}+3e^{-2\epsilon/kT}+5e^{-6\epsilon/kT}$$
Since the $j=2$ term yields a more negative numerator in the exponential, clearly the $j=1$ term matters more than it. Further terms will just have an even more negative value in the exponential, so we get the following approximation.
$$\implies Z_\text{rot}\approx 1+3e^{-2\epsilon/kT}$$
With $\beta=1/kT$, this implies that
$$\implies Z_\text{rot}\approx 1+3e^{-2\beta\epsilon}$$\\~\\

We want the average energy of the system, which is obtained by the following equation.
\begin{equation}
    \overline{E}=-\partDeriv{}{\beta}\ln(Z)
\end{equation}
If we say that there is just one molecule for now, then the following is the heat capacity (or if many particles, consider it the average heat capacity).
\begin{equation}
    C=\partDeriv{\overline{E}}{T}
\end{equation}
Clearly, we'll start by computing $\overline{E}$.\\~\\

$$\overline{E}=-\partDeriv{}{\beta}\ln(Z)$$
$$\implies\overline{E}=-\frac{1}{Z}\partDeriv{Z}{\beta}$$
$$\implies\overline{E}=-\frac{1}{1+3e^{-2\beta\epsilon}}\partDeriv{}{\beta}(1+3e^{-2\beta\epsilon})$$
$$\implies\overline{E}=-\frac{1}{1+3e^{-2\beta\epsilon}}(-6\epsilon e^{-2\beta\epsilon})$$
$$\implies\overline{E}=\frac{6\epsilon e^{-2\beta\epsilon}}{1+3e^{-2\beta\epsilon}}$$
As $\beta\rightarrow\infty$ due to $T\rightarrow0$, this implies that the exponential in the denominator will begin to be insignificant compared to the 1 added to it. Therefore this equation can be simplified to the following.
$$\implies\overline{E}=\frac{6\epsilon e^{-2\beta\epsilon}}{1}$$
$$\implies\overline{E}=6\epsilon e^{-2\beta\epsilon}$$
$$\implies\overline{E}=6\epsilon e^{-2\epsilon/kT}$$\\~\\

Now that we have an equation for $\overline{E}$, we can get the heat capacity $C$.
$$C=\partDeriv{\overline{E}}{T}$$
$$\implies C=\partDeriv{}{T}(6\epsilon e^{-2\epsilon/kT})$$
$$\implies C=6\epsilon \partDeriv{}{T}(e^{-2\epsilon/kT})$$
$$\implies C=6\epsilon e^{-2\epsilon/kT}\partDeriv{}{T}(-\frac{2\epsilon}{kT})$$
$$\implies C=6\epsilon e^{-2\epsilon/kT}\frac{2\epsilon}{kT^2}$$
$$\implies C=\frac{12\epsilon^2}{kT^2}e^{-2\epsilon/kT}$$\\~\\

To validate that this follows the third law of thermodynamics ($T\rightarrow0\implies C\rightarrow0$), I'll apply the limit.
$$C(T\rightarrow0)=\lim_{T\rightarrow0}\left(\frac{12\epsilon^2}{kT^2}e^{-2\epsilon/kT}\right)$$
$$\implies C(T\rightarrow0)=\frac{12\epsilon^2}{k}\cdot\lim_{T\rightarrow0}\left(\frac{e^{-2\epsilon/kT}}{T^2}\right)$$
Exponentials decay faster than quadratics do (I don't want to try to figure out L'Hopital's rule for this), so it's safe to say that the third law of thermodynamics applies.



\end{flushleft}

\end{document}
